% Copyright (c) 2026 Marcio Luiz Lima de Oliveira and V. el Hakim.
% Distribution and modification are subject to LPPL version 1.3c.
% SPDX-License-Identifier: LPPL-1.3c
% Part of the LPPL-covered work listed in MANIFEST.txt.
% See LICENSE and LICENSING.txt for the licence and scope.

% Visual design retained from the source thesis. See Chapters 6 and 7.
%Configurations

%%%
\usepackage[final]{microtype}
\microtypesetup{expansion=false} %Coment this line
\usepackage[dvipsnames]{xcolor}
\usepackage{intcalc}
\usepackage{graphicx}
\usepackage{amsmath}
\usepackage{amsthm}
\usepackage{amssymb}
\usepackage{mathtools}
\usepackage{subfig}
\usepackage{varwidth}
\usepackage[printonlyused,withpage]{acronym} %nohyperlinks
\usepackage{tikz}
\usepackage[most]{tcolorbox}
\usepackage[siunitx]{circuitikz}
\usepackage{algorithm}
\usepackage{algorithmicx}
\usepackage{algpseudocode}
\usepackage{enumitem}
\usepackage{minibox}
\usepackage{multirow}
\usepackage{wrapfig,booktabs}
\usepackage{setspace}
\usepackage{multicol}
\usepackage{colortbl}
\usepackage[pagestyles]{titlesec}
\usepackage{geometry}
\usepackage[bottom]{footmisc}
\usepackage{longtable}
\usepackage{tabularx}
% Additional flexibility only for paragraphs that otherwise cannot fit long
% literal command names in this compact book format.
\emergencystretch=3em
%\usepackage[T1,T2A]{fontenc}
\usepackage[utf8]{inputenc}
\usepackage[T1]{fontenc}
\usepackage{lmodern}
\usepackage[dutch,english]{babel}
\usepackage{alphabeta}
\usepackage{csquotes}
\usepackage[backend=bibtex,style=alphabetic, labelnumber, 
defernumbers=true,backref=true,url=false,
doi=false, isbn=false]{biblatex} %mine

\usepackage{listings} %mine
\usepackage{pdfpages} %pdf pages
\usepackage{ifthen}
%\usepackage[yyyymmdd]{datetime}
\usepackage[pages=some]{background} 
% contour is unused by the active cover; omit its engine-specific driver.
%\usepackage{eso-pic}

% Load hyperref after math, language, bibliography and float packages so their
% commands are patched consistently. Keep hyperfootnotes disabled as before.
\usepackage[hyperfootnotes=false]{hyperref}

%% Libraries
\tcbuselibrary{theorems}
\tcbuselibrary{skins}
\usetikzlibrary{arrows}
\usetikzlibrary{positioning}
\usetikzlibrary{calc}
\usetikzlibrary{backgrounds}
\usetikzlibrary{decorations.pathreplacing}

%%%%%%%%%%%%%%%%%%%%%%%%%%%%%%%%%%%%%%%%%%%%%%%%%%%%%%%%
% Algorithm for autoref
\providecommand{\algorithmautorefname}{Algorithm}
% Creates a LineComment for the Algorithm environment
\algnewcommand{\LineComment}[1]{\State \(\triangleright\) #1}
%%%
% Page size
%%
%% Set the paper size to 17 by 24 cm, approximately halfway between A4 and A5.
%\if@print
%% In print mode we add a 3 mm bleed on the outside of the page.
\geometry{
	papersize = {176mm,246mm},
	layoutsize = {170mm,240mm},
	layoutoffset = {3mm,3mm},
	bindingoffset = -3mm,
	inner = 25mm,
	outer = 20mm,
	top = 25mm,
	bottom = 22mm
}
% 2026-10-21: synchronize memoir's physical stock with geometry's final paper
% size. memoir otherwise restores its default Letter PDF canvas at begin-document.
% This hook also honours later personal or propositions geometry overrides.
\AddToHook{begindocument/before}{%
  \setstocksize{\paperheight}{\paperwidth}%
}
%\else
%\geometry{papersize={170mm,240mm}}
%\fi
%% We decrease the margins slightly from the default (scale = 0.7).
%\geometry{hscale=0.75,vscale=0.8}

%Add pdf cover to the first page
\backgroundsetup{contents={}} % Covers are drawn explicitly in Cover/BackCover.tex.

%%%%%%%%%%%%%%%%%%%%%%%%%%%%%%%%%%%%%%%%%%%%%%%%%%%%%%%%
%Hyper Reference Setup

\hypersetup{
	bookmarksopen=true,
	hypertexnames=true, %hypertexnames=false,
	colorlinks=true,% Set to false to disable coloring links
	citecolor= black!25!MyDarkTeal, %magenta,% The color of citations
	linkcolor=MyRed,% The color of references to document elements (sections, figures, etc)
	urlcolor=Purple!80!black,%blue,% The color of hyperlinks (URLs)
	%pdfstartview={FitV},
	%unicode,
	%breaklinks=true,
}

% Acronym links remain black using the package's public font hooks.
% Preserve the original fonts, including any future smaller-font option.
\NewCommandCopy{\ThesisOriginalAcsFont}{\acsfont}
\NewCommandCopy{\ThesisOriginalAcfFont}{\acffont}
\NewCommandCopy{\ThesisOriginalAcfsFont}{\acfsfont}
\renewcommand*{\acsfont}[1]{%
  \begingroup\hypersetup{hidelinks}\ThesisOriginalAcsFont{#1}\endgroup}
\renewcommand*{\acffont}[1]{%
  \begingroup\hypersetup{hidelinks}\ThesisOriginalAcfFont{#1}\endgroup}
\renewcommand*{\acfsfont}[1]{%
  \begingroup\hypersetup{hidelinks}\ThesisOriginalAcfsFont{#1}\endgroup}

%Command for time
%\newcommand{\ifnumber}[1]{\ifnum #1<9 {0#1}\else {#1}\fi} %Puts a 0 in front of numbers bellow 10
%\newcommand{\savedate}{\number\year\_\ifnumber{\number\month}\_\ifnumber{\number\day}}
%\newcommand{\chapterlocation}[1]{Chapters/Chapter_#1/Chapter_0#1}

%%%%%%%%%%%%%%%%%%%%%%%%%%%%%%%%%%%%%%%%%%%%%%%%%%%%%%%%

%New title command

\newcommand{\MyThesis}[6][]{
	\newcommand{\thesistitleA}{#1}
    \newcommand{\thesistitleB}{#2}
	\newcommand{\thesissubtitle}{#3}
	\newcommand{\thesisauthor}{#4}
	\newcommand{\thesisauthorshort}{#5}
	\newcommand{\thesisyear}{#6}
}

\newcommand{\MyDefense}[4][]{
	\newcommand{\defenseweekday}{#1}
	\newcommand{\defenseday}{#2}
	\newcommand{\defensemonth}{#3}
	\newcommand{\defensetime}{#4}
}

%%%%%%%%%%%%%%%%%%%%%%%%%%%%%%%%%%%%%%%%%%%%%%%%%%%%%
%Marcio's commands:

\newcommand{\myfont}{\sffamily} %Change the font of the text
\renewcommand*\familydefault{\sfdefault} %Change the font of the Chapter, Section, etc

\newcommand{\minsize}{0.85cm}
\newcommand{\myfontsize}{\scriptsize}
\newcommand{\myfontconfig}{\myfont \bfseries}

\renewcommand{\sectionautorefname}{Section}  %Change \autoref{} from section to Section
\renewcommand{\subsectionautorefname}{Subsection}
\renewcommand{\chapterautorefname}{Chapter}

\newcommand{\FigLocation}{Chapters/Chapter_0/fig/}
\newcommand{\ChapLocation}{Chapters/Chapter_0/}

%%%%%%%%%%%%%%%%%%%%%%%%%%%%%%%%%%%%%%%%%%%%%%%%%%%%%%%%
%Colors

\definecolor{DarkTeal}{HTML}{119E94}
\definecolor{NewTeal}{HTML}{1EC6BA} 

\definecolor{MyTeal}{HTML}{009999}
\definecolor{MyDarkTeal}{HTML}{269C95}
\definecolor{MyRed}{HTML}{A50026}

\definecolor{MarineBlue}{HTML}{20628E} 
\definecolor{LightMarineBlue}{HTML}{3684B7} 

\definecolor{Green}{HTML}{158E19}
\definecolor{LightGreen}{HTML}{21A825}

\definecolor{Blue}{HTML}{07388c}%{003CFF} 
\definecolor{LightBlue}{HTML}{256EFF}%{256EFF} 

\definecolor{BlueishGreen}{HTML}{00CC66} 
\definecolor{LightBlueishGreen}{HTML}{00FF80} 

\definecolor{MusgGreen}{HTML}{66CC00} 
\definecolor{LightMusgGreen}{HTML}{80FF00} 

\definecolor{Purple}{HTML}{4C0DAB} 
\definecolor{LightPurple}{HTML}{7B35E4}

\definecolor{PurpleBlue}{HTML}{6666FF} 
\definecolor{LightPurpleBlue}{HTML}{9999FF} 

\definecolor{BlueSky}{HTML}{3399FF} 
\definecolor{LightBlueSky}{HTML}{99CCFF} 

\definecolor{GreyBlue}{HTML}{679eb5} 
\definecolor{LightGreyBlue}{HTML}{77a3b5} 


\definecolor{Red}{HTML}{C92020} 
\definecolor{LightRed}{HTML}{E55555} 

%new 2023
\definecolor{CadetBlue}{HTML}{5F9EA0}
\definecolor{CornflowerBlue}{HTML}{6495ED}
\definecolor{EgyptianBlue}{HTML}{1434A4}


%%

\newcommand{\ColorBlue}{white!15!Blue}
\newcommand{\ColorLightBlue}{LightBlue}
\newcommand{\ColorGreen}{Green}
\newcommand{\ColorLightGreen}{LightGreen}
\newcommand{\ColorDarkTeal}{black!10!DarkTeal}
\newcommand{\ColorNewTeal}{NewTeal}
\newcommand{\ColorMarineBlue}{MarineBlue}
\newcommand{\ColorLightMarineBlue}{LightMarineBlue}
\newcommand{\ColorRed}{Red}
\newcommand{\ColorLightRed}{LightRed}
\newcommand{\ColorBlueishGreen}{black!10!BlueishGreen}
\newcommand{\ColorLightBlueishGreen}{LightBlueishGreen}
\newcommand{\ColorMusgGreen}{black!15!MusgGreen}
\newcommand{\ColorLightMusgGreen}{LightMusgGreen}
\newcommand{\ColorPurple}{Purple}
\newcommand{\ColorLightPurple}{LightPurple}

\newcommand{\ColorPurpleBlue}{black!20!PurpleBlue}
\newcommand{\ColorLightPurpleBlue}{LightPurpleBlue}
\newcommand{\ColorBlueSky}{black!17!BlueSky}
\newcommand{\ColorLightBlueSky}{LightBlueSky}

\newcommand{\ColorGreyBlue}{black!5!GreyBlue}
\newcommand{\ColorLightGreyBlue}{LightGreyBlue}

\newcommand{\ColorMyTeal}{white!10!MyTeal}

%new 2023
\newcommand{\ColorCadetBlue}{black!5!CadetBlue}
\newcommand{\ColorLightCadetBlue}{white!88!CadetBlue}
\newcommand{\ColorCornflowerBlue}{black!10!CornflowerBlue}
\newcommand{\ColorLightCornflowerBlue}{white!93!CornflowerBlue}
\newcommand{\ColorEgyptianBlue}{white!15!EgyptianBlue}
\newcommand{\ColorLightEgyptianBlue}{white!92!EgyptianBlue}

%%%%%%%%%%%%%%%%%%%%%%%%%%%%%%%%%%%%%%%%%%%%%%%%%%%%%
% Other Colors
%%
\definecolor{PinkNude}{HTML}{F89D7D}
\definecolor{LightTeal1}{HTML}{3e8680}
\definecolor{LightTeal2}{HTML}{175650}
\colorlet{LightTeal3}{white!30!LightTeal1}
\newcommand{\HeadingColor}{LightTeal3}
\newcommand{\TitleColor}{black!30!LightTeal2}
\newcommand{\titlestyle}{\normalfont\LARGE\scshape}

%%%%%%%%%%%%%%%%%%%%%%%%%%%%%%%%%%%%%%%%%%%%%%%%%%%%%
%% Maths Color Boxes: Theorem, Propositions, Definitions

\newtcbtheorem[number within=section]{theorem}{Theorem}
{
	theorem style = plain apart, enhanced, frame hidden,
	boxrule=2mm,titlerule=0mm,toptitle=1mm,bottomtitle=1mm,
	fonttitle=\bfseries\large, fontupper=\normalsize,
	coltitle=LightTeal2, colbacktitle=white!80!LightTeal2,
	colback=white!90!LightTeal2
}{thrm}
\newtcbtheorem[number within=section]{proposition}{Proposition}
{
	theorem style = plain apart, enhanced, frame hidden,
	boxrule=2mm,titlerule=0mm,toptitle=1mm,bottomtitle=1mm,
	fonttitle=\bfseries\large, fontupper=\normalsize,
	coltitle=LightTeal2, colbacktitle=white!80!LightTeal2,
	colback=white!90!LightTeal2
}{prop}
\newtcbtheorem[number within=section]{definition}{Definition}
{
	theorem style = plain apart, enhanced, frame hidden, breakable,
	boxrule=2mm,titlerule=0mm,toptitle=1mm,bottomtitle=1mm,
	fonttitle=\bfseries\large, fontupper=\normalsize,
	coltitle=black!60!PinkNude, colbacktitle=white!70!PinkNude,
	colback=white!80!PinkNude
}{def}
%%%%%%%%%%%%%%%%%%%%%%%%%%%%%%%%%%%%%%%%%%%%%%%%%%%%%%%%
% Page Colors
%%
\newcounter{chapCnt}
\setcounter{chapCnt}{-1} %Count Chapters


\definecolor{ThumbEnd}{HTML}{f5cac3}
\newcommand{\ThumbColIdx}[1]{white!2#1!\HeadingColor}

%\newcommand\ThumbColor{%
%	\ifcase \thechapCnt \TitleColor!0!ThumbEnd%
%	                   \or \TitleColor!14.2857!ThumbEnd%
%	                   \or \TitleColor!28.5714!ThumbEnd%
%	                   \or \TitleColor!42.8571!ThumbEnd%
%	                   \or \TitleColor!57.1429!ThumbEnd%
%	                   \or \TitleColor!71.4286!ThumbEnd%
%	                   \or \TitleColor!85.7143!ThumbEnd%
%	                   \or \TitleColor!100!ThumbEnd%
%	                   \else \TitleColor \fi}
                   

%Change 25/07/23
                   
%\newcommand{\ThumbColor}[1]{%
%	\ifcase #1 \ColorDarkTeal %1 
%	\or \ColorMarineBlue %2
%	\or \ColorGreyBlue %3
%	\or \ColorBlueSky %4
%	\or \ColorBlue %5
%	\or \ColorPurpleBlue %4
%	%\or \ColorRed
%	\or \ColorGreen
%	\or \ColorPurple
%	\else \ColorMyTeal \fi}

%\newcommand{\absColor}[1]{%
%	\ifcase #1 \ColorNewTeal %3 
%	\or \ColorLightMarineBlue %2
%	\or \ColorLightGreyBlue %5
%	\or \ColorLightBlueSky %6
%	\or \ColorLightBlue %1
%	\or \ColorLightPurpleBlue %4
%	%\or \ColorLightRed
%	\or \ColorLightGreen
%	\or \ColorLightPurple
%	\else \ColorMyTeal \fi}

 \newcommand{\ThumbColor}[1]{%
	\ifcase #1 \ColorDarkTeal %1
    \or \ColorCadetBlue %2
    \or \ColorGreyBlue %3
    \or \ColorMarineBlue %4
	\or \ColorBlueSky %5
    \or \ColorCornflowerBlue %6
    \or \ColorEgyptianBlue %7
	\or \ColorBlue %8
	\or \ColorPurpleBlue 
	%\or \ColorRed
	\or \ColorGreen
	\or \ColorPurple
	\else \ColorMyTeal \fi}

\newcommand{\absColor}[1]{%
	\ifcase #1 \ColorNewTeal %1
    \or \ColorLightCadetBlue %2
    \or \ColorLightGreyBlue %3
    \or \ColorLightMarineBlue %4
	\or \ColorLightBlueSky %5
    \or \ColorLightCornflowerBlue %6
    \or \ColorLightEgyptianBlue %7
	\or \ColorLightBlue %8
	\or \ColorLightPurpleBlue 
	%\or \ColorLightRed
	\or \ColorLightGreen
	\or \ColorLightPurple
	\else \ColorMyTeal \fi}

%Marcio' comands
\newcommand{\ColorTheme}{\ThumbColor{\thechapCnt}}
% Chapter number is drawn before chapCnt advances; derive its color from that
% counter instead of maintaining a second counter that can drift out of sync.
\newcommand{\NumberChapterColor}{\ThumbColor{\numexpr\value{chapCnt}+1\relax}}

%%%%%%%%%%% THUMB - MARK ON THE SIDE

\newlength\thumbheight
\setlength\thumbheight{35pt}%mine %{75pt} original
\newlength\thumbedge
\setlength\thumbedge{9.5pt}%mine  %{2.5pt} original
\newlength\thumbhspace
\setlength\thumbhspace{40.0pt}
\newlength\thumbvspace
%\setlength\thumbvspace{2.25cm}
\newcommand{\ThesisChapterCount}{8} % Number of numbered chapters; see Chapter 6.
\setlength\thumbvspace{(\textheight - 1cm - \thumbheight/2)/\ThesisChapterCount}

%%Draw the thumb

\newcommand*\lthumb{%
	\begin{tikzpicture}[remember picture,overlay]
	\coordinate (top margin) at (0pt,1in+\topmargin+\headheight+\headsep);
	\coordinate (left margin) at (1in+\evensidemargin,0pt);
	%% Calculate the corners of the thumb index based on the current
	%% chapter number.
	\coordinate (top left) at ($(current page.north west)-(top margin)-(0pt,\thechapCnt\thumbvspace + 0.5cm)$);
	\coordinate (bottom right) at ($(top left)+(left margin)-(\thumbhspace,\thumbheight)$);
	%% Shift the left edge to prevent the rounded corner from showing.
	\coordinate (top left) at ($(top left)-(\thumbedge,0pt)$);
	%% Draw the thumb index.
	\fill[fill=\ColorTheme,rounded corners=\thumbedge](top left) rectangle (bottom right);
	\end{tikzpicture}%
}
\newcommand*\rthumb{%
	\begin{tikzpicture}[remember picture,overlay]
	\coordinate (top margin) at (0pt,1in+\topmargin+\headheight+\headsep);
	\coordinate (right margin) at (1in+\evensidemargin,0pt);
	%% Calculate the corners of the thumb index based on the current
	%% chapter number.
	\coordinate (top right) at ($(current page.north east)-(top margin)-(0pt,\thechapCnt\thumbvspace+0.5cm)$);
	\coordinate (bottom left) at ($(top right)-(right margin)-(-\thumbhspace,\thumbheight)$);
	%% Shift the left right to prevent the rounded corner from showing.
	\coordinate (top right) at ($(top right)+(\thumbedge,0pt)$);
	%% Draw the thumb index.
	\fill[fill=\ColorTheme,rounded corners=\thumbedge](top right) rectangle (bottom left);
	\end{tikzpicture}%
}


%%%%%%%%%%% HEADERS AND FOOT NOTES - TOP AND BOTTOM OF THE PAGE

%\makeoddfoot{page style name}{stuff in the left}{stuff in the middle}{stuff in the right}

%Marcio command
%Create the style of the page number
%Style 1
%\newcommand{\PageNumberStyle}{              
%	\node(O)[fill=\ColorTheme, circle,minimum size = \minsize] %circle
%	{\textcolor{white}{\myfontconfig \myfontsize\thepage}}
%	}
%Style 2

%%%%%%
%My Page Circle (For numbering)

%Black Color
%\newcommand{\PageNumberStyle}{              
%	\node(O)[draw=\ColorTheme,very thick, circle,minimum size = \minsize, fill=\absColor{\thechapCnt}!25!white] %circle
%	{\textcolor{black}{\myfontconfig \myfontsize\thepage}}
%}

%White Color
\newcommand{\PageNumberStyle}{              
	\node(O)[draw=\ColorTheme,very thick, circle,minimum size = \minsize, fill=\ColorTheme] %circle
	{\textcolor{white}{\myfontconfig \myfontsize\thepage}}
}
%%%%%%
\makepagestyle{mypagestyle}
%Number  in Chapter First Page
\makeoddhead{plain}{}{}{}
%PageNumber New Chaper
%\makeoddfoot{plain}{}{\sffamily\bfseries\textcolor{\TitleColor}{\thepage}}{} %Original
%%%Marcio
\makeoddfoot{plain}{}{}{\myfontconfig \textcolor{\TitleColor}{{\begin{tikzpicture}[baseline=(O.base)]
				\PageNumberStyle;
		\end{tikzpicture}}
}}


%%%%

%\makeheadrule{mypagestyle}{\textwidth}{0.5pt} %original
%\makeheadfootruleprefix{mypagestyle}{\color{\HeadingColor}\vspace{-0.5pt}}{\color{\HeadingColor}} % old - make line
\makeheadrule{mypagestyle}{\textwidth}{1.5pt}
\makeheadfootruleprefix{mypagestyle}{\color{\ColorTheme}\vspace{-0.1pt}}{\color{\ColorTheme}}

\makeevenhead{mypagestyle}{\textcolor{black}{\leftmark}}{}{} 
\makeevenfoot{mypagestyle}{\lthumb\hspace{-1mm}\begin{tikzpicture}[baseline=(O.base)]
	%\node(O)[fill=\HeadingColor, rounded corners = 2.5,minimum size = 0.85cm] %original square
	\PageNumberStyle; %Number in page
	\end{tikzpicture}}{}{}

\makeoddhead{mypagestyle}{}{}{\textcolor{black}{\rightmark}} 
\makeoddfoot{mypagestyle}{}{}{\begin{tikzpicture}[baseline=(O.base)]
	%\node(O)[fill=\HeadingColor, rounded corners = 2.5,minimum size = 0.85cm] %original square
	\PageNumberStyle;
	\end{tikzpicture}\hspace{-1mm}\rthumb}

%\renewcommand{\chaptermark}[1]{ \markboth{\sffamily\bfseries \thechapter. #1}{} } %Original
%\renewcommand{\sectionmark}[1]{ \markright{\sffamily\bfseries\thesection. #1}{} } %Original
\renewcommand{\chaptermark}[1]{ \markboth{\myfontconfig Chapter \thechapter: #1}{} }
\renewcommand{\sectionmark}[1]{ \markright{\myfontconfig Section \thesection : #1}{} }


%%%%%%%%%%% APENDIX CHAPTERS STYLE

\makepagestyle{notationstyle}
\makeheadrule{notationstyle}{\textwidth}{1.5pt}
\makeheadfootruleprefix{notationstyle}{\color{\ColorTheme}\vspace{-0.1pt}}{\color{\ColorTheme}}

\makeevenhead{notationstyle}{\leftmark}{}{} 
\makeevenfoot{notationstyle}{\hspace{-1mm}\begin{tikzpicture}[baseline=(O.base)]
	%\node(O)[fill=\HeadingColor, rounded corners = 2.5,minimum size = 0.85cm]  %original square
	\PageNumberStyle;
	\end{tikzpicture}}{}{}

\makeoddhead{notationstyle}{}{}{\leftmark} 

\makeoddfoot{notationstyle}{}{}{\begin{tikzpicture}[baseline=(O.base)]
	%\node(O)[fill=\HeadingColor, rounded corners = 2.5,minimum size = 0.85cm] %original square
	\PageNumberStyle;
	\end{tikzpicture}\hspace{-1mm}}

%%%%%%%%%%%%%% FOOTNOTES
% One fixed separator for ordinary and publication notes. Neither helper changes
% shared page-output formatting; publication notes remain unnumbered.
\setfootnoterule{3pt}{2in}{0.4pt}
\DeclareRobustCommand{\pubfootnote}[1]{%
  \begingroup
    \renewcommand{\thefootnote}{}%
    \footnotetext[0]{#1}% No mark and no change to the ordinary footnote counter.
  \endgroup
}
\DeclareRobustCommand{\myfootnote}[1]{\footnote{#1}}

%%%%%%%%%%%%%%%%%%%%%%%%%%%%%%%%%%%%%%%%%%%%%%%%%%%%%%%%
% Chapter's style
%%
\makechapterstyle{thesis}{
		
		\setlength{\beforechapskip}{10pt}
		\setlength{\afterchapskip}{1.5cm}
		\setlength{\midchapskip}{1cm}
	\renewcommand*{\printchaptername}{}
	\renewcommand*{\chapnumfont}{\normalfont \myfontconfig \huge}
	\renewcommand*{\printchapternum}{       %Print number in the first page of the chapter
		%\flushright %original at the right corner
		\centering %Center without opening an unclosed center environment
		%Beginning of chapter
		\begin{tikzpicture}
		%\draw[fill,color=\HeadingColor,rounded corners = 5] (0,0) rectangle (2cm,2cm); %original rectangle
		
		%Mine
		\draw[fill,color=\NumberChapterColor,rounded corners = 15] (0,0) rectangle (2cm,2cm); 
		\draw[color=white] (1cm,1cm) node {\chapnumfont \thechapter }; %original
		
		\end{tikzpicture}
		\stepcounter{chapCnt}
	}
	%\renewcommand*{\chaptitlefont}{\normalfont \myfontconfig \Huge \color{\TitleColor}} %original
	\renewcommand*{\chaptitlefont}{\normalfont \myfontconfig \Huge \color{\ColorTheme!50!black}}
	%\renewcommand*{\printchaptertitle}[1]{\flushright\chaptitlefont##1} %original - chapter at right
	\renewcommand*{\printchaptertitle}[1]{\centering \chaptitlefont##1} %Chapter title in the middle
}
\chapterstyle{thesis}



%%%%%%%%%%%%%%%%%%%%%%%%%%%%%%%%%%%%%%%%%%%%%%%%%%%%%%%%
% Chapter's abstract style
%%
%theorem style = plain apart, enhanced, frame hidden,
%boxrule=2mm,titlerule=0mm,toptitle=1mm,bottomtitle=1mm,
%fonttitle=\bfseries\large, fontupper=\normalsize,
%coltitle=LightTeal2, colbacktitle=white!80!LightTeal2,
%colback=white!90!LightTeal2

%%%%%%%%%%%%% CHAPTER ABSTRACT:
\newtcolorbox{absbox}
{
	frame hidden,
	colback= white!90!\absColor{\thechapCnt}, %white!90!LightTeal2,
	fonttitle=\scshape, coltitle=\ColorTheme,%LightTeal2,
	fontupper=\small,
	before upper={\parindent15pt},
	colbacktitle=white!85!\absColor{\thechapCnt},
	enhanced,
	attach boxed title to top center={yshift=-2mm},title=Abstract,
	boxed title style= {arc=0pt, outer arc=0pt, boxrule=0pt, after skip = 1cm}
}

\abstractrunin % Places the abstract word before the text else it is place above and aligned with \absnamepos
% 'Abstract' same formatting as float caption
\setlength{\abstitleskip}{0pt}
\renewcommand{\abstractnamefont}{\normalfont\small\scshape}
\abslabeldelim{\;--\;}
%\setlength{\absleftindent}{\specialparindent}
%\setlength{\absrightindent}{\specialparindent}
\renewcommand{\abstracttextfont}{\normalfont\small\itshape}
\let\oldendabstract\endabstract
\def\endabstract{%
	\oldendabstract
	\fancybreak{\rule{0pt}{\onelineskip}}%
}

%%%%%%%%%%%%%%%%%%%%%%%%%%%%%%%%%%%%%%%%%%%%%%%%%%%%%%%%
% Fonts
%%
%\newcommand{\SectionColor}{\TitleColor}
\newcommand{\SectionColor}{\ColorTheme!30!black}

\titleformat{\section}
%original - with extra dot
%{\normalfont\fontsize{13}{0}\sffamily\bfseries\color{\TitleColor}}{\textcolor{\TitleColor}{\thesection.}} 
{\normalfont\fontsize{13}{16}\myfontconfig \color{\SectionColor}}{\textcolor{\SectionColor}{\thesection}} 
{1em}{} 
\titleformat{\subsection}
%original - with extra dot
%{\normalfont\fontsize{11}{0}\sffamily\bfseries\color{\TitleColor}}{\textcolor{\TitleColor}{\thesubsection.}}
{\normalfont\fontsize{11}{14}\myfontconfig \color{\SectionColor}}{\textcolor{\SectionColor}{\thesubsection}}
{1em}{} 
\titleformat{\subsubsection} {\normalfont \myfontconfig \color{\SectionColor}}{}{1em}{}




% TEMPLATE CONTROLS. Chapter 6 explains these and the original design above.
% Keep the original page/heading/theme helpers; hide only the edge markers.
\newif\ifThesisThumbs
\ThesisThumbstrue
\let\ThesisLeftThumb\lthumb
\let\ThesisRightThumb\rthumb
\renewcommand{\lthumb}{\ifThesisThumbs\ThesisLeftThumb\fi}
\renewcommand{\rthumb}{\ifThesisThumbs\ThesisRightThumb\fi}
% Source examples are literal: macros inside listings are never executed.
\lstset{language=[LaTeX]TeX,basicstyle=\ttfamily\footnotesize,
  columns=fullflexible,keepspaces=true,breaklines=true,showstringspaces=false,
  frame=single,rulecolor=\color{black!15},backgroundcolor=\color{black!2},
  xleftmargin=0.5em,xrightmargin=0.5em,aboveskip=1em,belowskip=1em}
